\chapter{Supersymmetric Wilsonian Schemes and Exact Dynamics}
\label{ch:susy-wilsonian-schemes-exact-dynamics}

\ConventionsLorentzian

Exact statements in supersymmetric QFT require a declared Wilsonian scheme.
The scheme is the data that define which modes are integrated out, which
local coordinates are used for the resulting action, which Ward identities
hold at finite cutoff, and which anomaly classes obstruct stronger
statements.  Supersymmetry of the Wilsonian integration and supersymmetry of
the regulator are separate requirements.

\section{Wilsonian Data}

\begin{definition}[Supersymmetric Wilsonian datum]
\label{def:susy-wilsonian-datum}
A supersymmetric Wilsonian datum at scale \(\Lambda\) consists of:
a space \(\mathcal E_\Lambda\) of regulated superfield variables, a gauge
group action together with either a BV gauge complex or a lattice
gauge-field construction, a projection
\(\pi_{\Lambda\Lambda'}:\mathcal E_{\Lambda}\to\mathcal E_{\Lambda'}\) for
\(\Lambda>\Lambda'\), an integration cycle in the discarded variables, and a
local action functional \(S_\Lambda\) satisfying the declared gauge and
supersymmetry Ward identities at finite cutoff.
\end{definition}

If the gauge symmetry is present, the natural finite-cutoff statement is a BV
statement.  Let \((\,\cdot\,,\,\cdot\,)_\Lambda\) be the odd BV bracket and
\(\Delta_\Lambda\) the regulated BV Laplacian.  The quantum master equation is
\[
  \frac12(S_\Lambda,S_\Lambda)_\Lambda
  -\ii\hbar\,\Delta_\Lambda S_\Lambda=0 .
\]
This equation is meaningful only after the regulator and the functional
measure entering \(\Delta_\Lambda\) have been specified.
For a gauge-fixed supersymmetric Wilsonian path integral, gauge consistency is
the same equation before restriction to the gauge-fixing Lagrangian
submanifold.  The gauge-fixed Slavnov--Taylor identities and supersymmetry
Ward identities are obtained by restricting the BV half-density and applying
BV Stokes.  A Wilsonian coarse-graining step is gauge-consistent when it is a
BV pushforward that sends a solution of the quantum master equation at
\(\Lambda\) to a solution at \(\Lambda'\).

A lattice supersymmetric gauge construction, when it exists, is the other
finite-cutoff route: gauge fields are link variables, gauge invariance is
exact at finite lattice spacing, and Wilsonian coarse graining is a blocking
map for link variables or gauge-invariant observables.  The obstruction is
then shifted to the realization of supersymmetry, fermions, auxiliary
variables, reflection positivity, and the continuum limit.  The two routes are
therefore complementary: BV is the continuum homological framework for
gauge-compatible effective actions, while the lattice gives a compact
gauge-invariant finite integral.

\section{Coarse Graining as BV Pushforward}

For \(\Lambda>\Lambda'\), split the regulated variables into retained
variables \(\phi'\in\mathcal E_{\Lambda'}\) and eliminated variables
\(\chi\).  The Wilsonian action at \(\Lambda'\) is defined by
\[
  \exp\!\left(\frac{\ii}{\hbar}S_{\Lambda'}(\phi')\right)
  =
  \int_{\mathcal L_{\Lambda\Lambda'}}
  \exp\!\left(\frac{\ii}{\hbar}S_{\Lambda}(\phi',\chi)\right)D\chi ,
\]
where \(\mathcal L_{\Lambda\Lambda'}\) is the declared integration cycle.
When the split is compatible with the BV odd symplectic form, this is a BV
pushforward.  The basic theorem is formal but precise: BV pushforward sends
solutions of the quantum master equation to solutions, provided the boundary
term on \(\mathcal L_{\Lambda\Lambda'}\) vanishes.

For gauge theories, this BV compatibility or a lattice blocking construction
is part of the definition of the Wilsonian step.  A projection of gauge
potentials by itself does not define a gauge-theory coarse graining; the
projection must be lifted to the full ghost-antifield complex with a
compatible half-density and integration cycle, or replaced by a finite
gauge-invariant lattice blocking map.

\begin{proposition}[Semigroup property]
\label{prop:susy-wilsonian-semigroup}
Assume the integration cycles for
\(\Lambda>\Lambda'>\Lambda''\) compose by Fubini's theorem and have no
boundary contribution.  Then Wilsonian integration satisfies
\[
  S_{\Lambda''}
  =
  \mathcal R_{\Lambda'\Lambda''}
  \bigl(\mathcal R_{\Lambda\Lambda'}(S_\Lambda)\bigr)
  =
  \mathcal R_{\Lambda\Lambda''}(S_\Lambda),
\]
where \(\mathcal R\) denotes the corresponding Wilsonian map.
\end{proposition}

\begin{proof}
Insert the definition of \(S_{\Lambda'}\) into the definition of
\(S_{\Lambda''}\).  The iterated integral is the integral over the product
cycle of the two eliminated variable sets.  By the hypothesis on the cycles
and measure, this product cycle is the cycle defining direct integration from
\(\Lambda\) to \(\Lambda''\).  Taking the logarithm of the resulting
exponential functional gives the stated identity.
\end{proof}

\section{Local Coordinates: F-Terms and D-Terms}

For four-dimensional \(\mathcal N=1\) theories, local supersymmetric
functionals split by superspace measure:
\[
  \int d^4x\,d^2\theta\,\mathcal F(\Phi,W_\alpha,\nabla)
  +\int d^4x\,d^4\theta\,\mathcal K(\Phi,\bar\Phi,V,\nabla).
\]
The first class is called \(F\)-term coordinates and the second class
\(D\)-term coordinates.  This coordinate decomposition becomes a statement
about dynamics only after a Wilsonian scheme and its Ward identities have
been specified.  The gauge kinetic coordinate is the chiral functional
\[
  \frac{1}{16\pi\ii}\int d^4x\,d^2\theta\,
  \tau\,\operatorname{tr}(W^\alpha W_\alpha)+\text{c.c.},
  \qquad
  \tau=\frac{\theta}{2\pi}+\frac{4\pi\ii}{g^2},
\]
with the trace normalization fixed in
Chapter~\ref{ch:susy-gauge-theory}.  If a different superspace normalization
for \(W_\alpha\) is chosen, the displayed coefficient is changed by the
component-matching factor; the physical convention is the component
Yang--Mills term \(-\frac{1}{4g^2}\operatorname{tr}F^2\).

\section{Holomorphic Wilsonian Claims}

\begin{hypothesis}[Holomorphic Wilsonian scheme]
\label{hyp:holomorphic-wilsonian-scheme}
A Wilsonian scheme is holomorphic for a set of chiral couplings \(t^A\) if
the regulated chiral action depends holomorphically on \(t^A\), the
coarse-graining map preserves the chiral superspace cohomology class of local
\(F\)-terms, and all violations of holomorphy are represented by declared
local anomaly functionals.
\end{hypothesis}

Under this hypothesis, nonrenormalization statements become statements about
coordinates.  A superpotential coordinate
\[
  \int d^4x\,d^2\theta\, W_{\rm tree}(\Phi;t)
\]
can receive only those Wilsonian corrections that are holomorphic functions
of the chiral coordinates, respect the exact symmetries including anomalous
spurion transformations, and are local at scale \(\Lambda\).  The
selection-rule arguments are replaced here by these three stated conditions.

\begin{proposition}[Selection rule for a Wilsonian \(F\)-term]
\label{prop:susy-f-term-selection-rule}
Assume Hypothesis~\ref{hyp:holomorphic-wilsonian-scheme}.  Let
\(\mathcal U(t,\Phi)\) be a candidate correction to a chiral \(F\)-term.  If
\(\mathcal U\) is not a holomorphic section of the same symmetry line bundle
as the original chiral measure and operator insertion, then it is absent from
the Wilsonian \(F\)-term coordinate in that scheme.
\end{proposition}

\begin{proof}
The Wilsonian \(F\)-term coordinate is, by hypothesis, a holomorphic element
of the chiral local cohomology with specified transformation under exact and
anomalous symmetries.  Local terms with different transformation law lie in a
different line bundle over the coupling space.  They cannot be added to the
coordinate without violating the Ward identity defining the scheme.  Thus
only holomorphic sections with the same transformation law are admissible.
\end{proof}

\section{Exact Dynamics as a Scheme-Dependent Statement}

The chiral ring at scale \(\Lambda\) is the cohomology of local chiral
operators modulo \(\bar Q\)-exact operators and equations of motion in the
Wilsonian scheme.  A moduli space statement is an assertion about expectation
values of this ring in supersymmetric states together with the effective
action governing fluctuations.  A duality statement is an isomorphism between
two such Wilsonian data, including line operators, anomalies, contact terms,
and relevant deformations.

\begin{openproblem}[Wilsonian construction]
\label{op:susy-wilsonian-construction}
Construct a Wilsonian scheme for four-dimensional gauge theories that
preserves gauge invariance, realizes supersymmetry in BV form at finite
cutoff, states the anomaly functionals, and proves the holomorphic hypotheses
used in exact dynamics.  This is the infrastructure needed for proofs of
\(\mathcal N=1\) and \(\mathcal N=2\) gauge dynamics.
\end{openproblem}
